\documentclass[11pt, a4paper]{article}
\usepackage{graphicx} % Required for inserting images

% Configuración de idioma y codificación
\usepackage[T1]{fontenc}
\usepackage[utf8]{inputenc}
% Mejorar compatibilidad con captions y hooks
\usepackage{caption}
\usepackage[spanish]{babel}
\usepackage{csquotes}

% Configuración de bibliografía
\usepackage[backend=biber, style=apa]{biblatex}
\addbibresource{referencias.bib}

% Paquetes para matemáticas e imágenes
\usepackage{amsmath}
\usepackage{amssymb}

\title{Documento LLMs, VLLMs}   
\author{Facundo Gaviola}
\date{\today}

\begin{document}

\maketitle


Este documento explora la evolución técnica en el campo del Procesamiento de Lenguaje Natural (NLP), analizando la transición desde las arquitecturas secuenciales basadas
en Redes Neuronales Recurrentes (RNN) y Long Short-Term Memory (LSTM) hasta el cambio de paradigma introducido por los Transformers y posteriores LLMs y su evolucion 
hacias los VLLMs.


\section{¿Qué es el Procesamiento de Lenguaje Natural (NLP)?}

El Procesamiento de Lenguaje Natural (NLP, por sus siglas en inglés) es una disciplina fundamental dentro de la Inteligencia Artificial (IA) que ha evolucionado rápidamente
en los últimos años. Su objetivo central es capacitar a las máquinas para comprender, interpretar y generar lenguaje humano de manera eficiente.

Según Atkinson-Abutridy, el NLP es el corazón de las aplicaciones innovadoras actuales, como los sistemas conversacionales (chatbots). La esencia de los avances modernos en
NLP radica en la capacidad de capturar patrones complejos y aprender representaciones contextuales del lenguaje, superando las limitaciones de los antiguos sistemas basados
en reglas fijas.

En la era actual, el NLP se centra predominantemente en el uso de \textbf{Grandes Modelos de Lenguaje (LLMs)}, los cuales son redes neuronales con miles de millones de
parámetros entrenados con cantidades masivas de texto no etiquetado mediante aprendizaje autosupervisado.

\subsection{Utilidades y Aplicaciones del NLP}

El libro destaca que los modelos de lenguaje modernos permiten a las computadoras realizar una amplia gama de tareas. Las aplicaciones principales descritas incluyen:

\begin{enumerate}
    \item \textbf{Sistemas Conversacionales y Chatbots:} Como ChatGPT o BARD, capaces de mantener diálogos coherentes, responder preguntas de seguimiento y recordar detalles 
    previos de la conversación.
    \item \textbf{Clasificación de Sentimientos:} Análisis de opiniones para determinar la polaridad emocional (positiva, negativa o neutra) en textos provenientes de redes 
    sociales o reseñas de productos.
    \item \textbf{Generación de Resúmenes:} Capacidad para sintetizar documentos largos en versiones concisas manteniendo la información clave.
    \item \textbf{Traducción Automática:} Traducción de texto entre diferentes idiomas con alta precisión contextual.
    \item \textbf{Búsqueda Semántica:} A diferencia de la búsqueda por palabras clave, los modelos de NLP modernos permiten buscar basándose en el significado y la intención 
    detrás de la consulta del usuario.
    \item \textbf{Generación de Código:} Herramientas que traducen descripciones en lenguaje natural a código de programación funcional (e.g., Python, SQL).
    \item \textbf{Razonamiento y QA (Question Answering):} Sistemas capaces de responder preguntas abiertas o cerradas basándose en un contexto dado o en conocimiento general.
\end{enumerate}

\subsection{Taxonomía de Modelos}

En el aprendizaje automático aplicado al procesamiento del lenguaje y la visión, los modelos pueden clasificarse según el tipo
de distribución de probabilidad que intentan aprender y el objetivo final de su entrenamiento:

\begin{itemize}
    \item \textbf{Modelos Discriminativos:} Su objetivo principal es aprender la frontera de decisión entre clases. Modelan la
    probabilidad condicional $P(Y|X)$, donde $X$ es la entrada (e.g., una imagen o texto) e $Y$ es la etiqueta. Históricamente,
    estos modelos han dominado tareas como la clasificación de imágenes, donde arquitecturas como ResNet se entrenan de manera
    supervisada para predecir categorías fijas \parencite{radford_learning_nodate}. En NLP, modelos como BERT, aunque pre-entrenados, se
    utilizan frecuentemente en configuraciones discriminativas mediante el ajuste fino para clasificación \parencite{devlin_bert_nodate}.
    
    \item \textbf{Modelos Generativos:} A diferencia de los discriminativos, estos modelos aprenden la distribución conjunta
    $P(X,Y)$ o simplemente la distribución de los datos $P(X)$ si no hay etiquetas. Su objetivo es entender cómo se generan los
    datos para poder crear nuevas muestras que se asemejen a la distribución original. Los LLMs modernos, como GPT-3, pertenecen
    a esta categoría, ya que se entrenan para modelar la probabilidad de secuencias de texto \parencite{brown_language_nodate}.
\end{itemize}

\subsection{Modelos Generativos Profundos}

Dentro del aprendizaje profundo, los modelos generativos han evolucionado en diversas arquitecturas capaces de sintetizar datos
complejos como imágenes y texto. Podemos categorizarlos en cuatro familias principales según su mecanismo de aprendizaje:

\begin{enumerate}
    \item \textbf{Redes Generativas Antagónicas (GANs):} Se basan en un juego de suma
    cero entre dos redes: un \textit{generador} que crea muestras falsas y un \textit{discriminador} que intenta distinguir entre
    muestras reales y falsas. Aunque exitosas en generación de imágenes, su entrenamiento es notoriamente inestable 
    \parencite{goodfellow_generative_nodate}.
    
    \item \textbf{Modelos Basados en Flujo (Flow-based):}  Utilizan una secuencia de transformaciones
    invertibles para mapear una distribución simple a una compleja, permitiendo el cálculo exacto de la verosimilitud de los datos 
    \parencite{rezende_variational_nodate}.
    
    \item \textbf{Modelos de Difusión (Diffusion Models):} Aprenden a revertir un proceso de 
    difusión gradual que añade ruido a los datos hasta destruirlos. Al aprender a eliminar este ruido paso a paso, pueden generar 
    datos de alta calidad a partir de ruido aleatorio \parencite{ho_denoising_nodate}.
    
    \item \textbf{Modelos Autoregresivos:} Esta es la arquitectura fundamental detrás de los LLMs. Estos modelos descomponen la
    distribución de probabilidad de una secuencia en el producto de probabilidades condicionales:
    \begin{equation}
        P(x) = \prod_{i=1}^{n} P(x_i | x_1, ..., x_{i-1})
    \end{equation}
    En este enfoque, el modelo predice el siguiente elemento (token) de la secuencia basándose únicamente en los anteriores.
    Arquitecturas como el Transformer original en su fase de decodificación \parencite{vaswani_attention_2023} y modelos masivos como 
    GPT-3 operan bajo este principio, consumiendo los símbolos generados previamente como entrada adicional para generar el 
    siguiente \parencite{vaswani_attention_2023, brown_language_nodate}. Incluso en visión computacional, modelos como ImageGPT han aplicado este
    principio a secuencias de píxeles \parencite{dosovitskiy_image_2021}.
\end{enumerate}

%\subsection{Tipos de Modelos en NLP}

%La evolución de los modelos en NLP ha pasado desde enfoques estadísticos hasta redes neuronales profundas y arquitecturas basadas en Transformers. A grandes rasgos, los modelos
%se pueden caracterizar por su objetivo (generativos vs. discriminativos), por su arquitectura (encoders, decoders, o ambos) y por el tipo de tarea para la cual se optimizan 
%(clasificación, generación, traducción, etc.).

%Más allá de la distinción teórica, es útil considerar tres familias prácticas:
%\begin{enumerate}
%\item \textbf{Modelos basados en codificador (encoder-only):} están optimizados para comprender y representar texto. Estos modelos producen una representación contextual para
%cada token y se usan principalmente en tareas de comprensión como clasificación, extracción de información y respuesta a preguntas. BERT y sus derivaciones son ejemplos clásicos.
%\item \textbf{Modelos autorregresivos (decoder-only):} se entrenan para modelar la probabilidad condicional del siguiente token dado el contexto previo, lo que los hace especialmente
%buenos para generación de texto y completado. GPT y sus sucesores encajan en este grupo.
%\item \textbf{Modelos encoder-decoder (seq2seq):} combinan un encoder que resume la entrada y un decoder que genera la salida; son la elección natural para tareas de traducción 
%automática, summarización y tareas de transformación de texto a texto. T5 y la mayoría de los sistemas de traducción neuronal usan este enfoque.
%\end{enumerate}
%Cada familia tiene ventajas y limitaciones: los encoders capturan contexto bidireccional y son fuertes en comprensión, los decoders autorregresivos resultan flexibles para generación
%libre y continuaciones, y los encoder-decoder equilibran comprensión y generación para transformar una pieza de texto en otra.

%\subsubsection{ Modelos Generativos vs. Discriminativos}
%Es fundamental distinguir entre estas dos categorías:
%\begin{itemize}
%    \item \textbf{Modelos Discriminativos:} Se centran en clasificar datos o predecir una etiqueta $Y$ dada una entrada $X$ (calculan la probabilidad condicional $P(Y|X)$). 
%    Ejemplo: Decidir si un correo es spam o no.
%    \item \textbf{Modelos Generativos:} Aprenden la distribución de los datos y cómo se generan (calculan la probabilidad conjunta $P(X,Y)$). Tienen la capacidad de crear 
%    nuevas instancias de datos (texto, imágenes) que se asemejan a los datos reales. Los LLMs modernos son fundamentalmente generativos.
%\end{itemize}

\section{Representación del Lenguaje: Embeddings}

Para que los algoritmos de aprendizaje profundo puedan procesar el lenguaje natural, es necesario transformar las unidades léxicas (palabras o subpalabras) en representaciones 
numéricas operables. La evolución de estas representaciones ha sido clave para el éxito de los modelos modernos.

\subsection{One-Hot Encoding}
Históricamente, la aproximación más inmediata para representar texto fue la codificación \textit{One-Hot}. En este esquema, cada palabra del vocabulario $V$ se representa como
un vector binario disperso $v \in \mathbb{R}^{|V|}$, donde todos los elementos son cero excepto el índice correspondiente a la palabra en cuestión, que toma el valor de 1.

Aunque intuitivo, este método presenta dos limitaciones fundamentales para el modelado del lenguaje:
\begin{itemize}
    \item \textbf{Alta dimensionalidad y dispersión:} El tamaño del vector crece linealmente con el tamaño del vocabulario, lo que resulta computacionalmente ineficiente para
    corpus grandes.
    \item \textbf{Ortogonalidad semántica:} Al ser vectores ortogonales entre sí, el producto punto entre dos palabras cualesquiera es siempre cero. Esto impide capturar 
    relaciones semánticas; por ejemplo, el modelo no puede inferir matemáticamente que ''perro'' y ''gato'' son más similares entre sí que ''perro'' y ''automóvil''.
\end{itemize}

\subsection{Word Embeddings}
Para superar las limitaciones de las representaciones dispersas, surgió el concepto de \textit{Word Embeddings} o incrustaciones de palabras. Esta técnica consiste en mapear
palabras a vectores densos de valores reales en un espacio vectorial de menor dimensión ($d \ll |V|$), donde la similitud semántica entre palabras se traduce en proximidad
geométrica \parencite{atkinson2025}.

Un hito fundacional en esta área fue la propuesta de \textcite{mikolov_efficient_2013} con la arquitectura conocida popularmente como \textbf{Word2Vec}. Los autores introdujeron técnicas
eficientes para estimar estas representaciones a partir de grandes corpus de texto no estructurado, basándose en la hipótesis distribucional: palabras que aparecen en contextos
similares tienden a tener significados similares.

\textcite{mikolov_efficient_2013} propusieron dos arquitecturas principales de redes neuronales superficiales para aprender estos pesos:
\begin{enumerate}
    \item \textbf{Continuous Bag-of-Words (CBOW):} El modelo predice la palabra actual basándose en el promedio de los vectores de las palabras de su contexto (pasado y futuro).
    \item \textbf{Continuous Skip-gram:} El modelo utiliza la palabra actual para predecir las palabras que aparecen en su contexto circundante dentro de una ventana definida.
\end{enumerate}

La innovación clave de estos modelos fue su capacidad para capturar relaciones sintácticas y semánticas complejas mediante operaciones algebraicas simples sobre los vectores. Como demostraron
\textcite{mikolov_efficient_2013}, el resultado de la operación vectorial $vec(\text{Rey}) - vec(\text{Hombre}) + vec(\text{Mujer})$ resulta en un vector muy cercano a $vec(\text{Reina})$. 
Este avance demostró que las redes neuronales podían aprender representaciones distribuidas del lenguaje de alta calidad con un costo computacional significativamente menor que los 
enfoques anteriores.

\section{Arquitecturas Secuenciales}

\subsection{Fundamentos: El Perceptrón Multicapa y la Aproximación Universal}

Antes de abordar el procesamiento de secuencias, es vital comprender la capacidad teórica de las redes neuronales precursoras. Las Redes Neuronales (NNs) tradicionales, o Perceptrones
Multicapa (MLP), operan bajo una arquitectura \textit{feed-forward}, donde la información fluye unidireccionalmente desde la entrada hasta la salida.

El éxito de estas arquitecturas no es accidental. En un trabajo seminal, \textcite{hornik_multilayer_1989} demostraron rigurosamente que las redes \textit{feed-forward} multicapa son una clase de
\textbf{aproximadores universales}. Su teorema establece que una red con al menos una capa oculta y suficiente cantidad de neuronas puede aproximar cualquier función Borel medible (es decir, 
prácticamente cualquier función continua de interés práctico) con un grado de error arbitrariamente pequeño, siempre que la función de activación sea una ''squashing function'' (función acotada y
no decreciente, como la sigmoide o tangente hiperbólica) \parencite{hornik_multilayer_1989}.

Este resultado garantiza teóricamente que una red neuronal tiene la capacidad estructural para modelar relaciones complejas. Sin embargo, estas redes poseen una limitación fundamental para el 
lenguaje: asumen que cada entrada es independiente de la anterior, careciendo de memoria para capturar el contexto secuencial.

\subsection{Redes Neuronales Recurrentes (RNNs)}

Para superar la falta de memoria de los MLPs, surgieron las Redes Neuronales Recurrentes (RNNs). A diferencia de las redes tradicionales, las RNNs introducen ciclos de retroalimentación y el concepto
de \textbf{pesos compartidos} a través del tiempo. Esto permite que la red procese una secuencia no como entradas aisladas, sino manteniendo un estado interno que evoluciona paso a paso.

Un hito fundamental en esta evolución fue el trabajo de \textcite{elman_finding_1990}, titulado \textit{''Finding Structure in Time''}. Elman revolucionó el enfoque argumentando que el tiempo no debía
representarse espacialmente (como una secuencia de entradas simultáneas en una ventana fija), sino de manera \textbf{implícita} a través de su efecto en el procesamiento interno. Para ello, introdujo 
la arquitectura conocida como \textbf{Red Recurrente Simple (SRN)}.

La innovación clave de la SRN fue la introducción de \textbf{unidades de contexto}, las cuales almacenan una copia exacta de las activaciones de la capa oculta en el paso de tiempo anterior ($h_{t-1}$). 
Así, la red procesa la entrada actual combinándola con su propia ''historia'' interna. Matemáticamente, el estado oculto $h_t$ se define por:

\begin{equation}
    h_t = f(W_{in} \cdot x_t + W_{ctx} \cdot h_{t-1} + b)
\end{equation}

Donde:
\begin{itemize}
    \item $h_t$: Vector de estado oculto en el tiempo $t$.
    \item $f$: Función de activación no lineal (usualmente sigmoide o tanh).
    \item $W_{in}$: Matriz de pesos que conecta la entrada con la capa oculta.
    \item $x_t$: Vector de entrada en el tiempo $t$.
    \item $W_{ctx}$: Matriz de pesos recurrentes (o de contexto) que conecta el estado pasado con el actual.
    \item $h_{t-1}$: Estado oculto del paso anterior (copiado en las unidades de contexto).
    \item $b$: Vector de sesgo (bias).
\end{itemize}

En su estudio, Elman demostró que estas redes podían:
\begin{itemize}
    \item Aprender lógica temporal compleja (como la función XOR secuencial).
    \item Descubrir la noción de ''palabra'' en un flujo continuo de letras, utilizando el error de predicción como detector de límites.
    \item Desarrollar representaciones internas jerárquicas que agrupaban palabras en categorías gramaticales (e.g., verbos vs. sustantivos) basándose únicamente en la estadística de su aparición 
    secuencial \parencite{elman_finding_1990}.
\end{itemize}

Sin embargo, a pesar de estos avances en la representación estructural, las RNNs simples como la de Elman enfrentan una limitación crítica conocida como el problema del 
\textit{\textbf{desvanecimiento del gradiente}} (vanishing gradient). A medida que la secuencia se alarga, los gradientes utilizados para actualizar los pesos se vuelven infinitesimales, impidiendo
que la red aprenda dependencias a muy largo plazo \parencite{atkinson2025}. Esta limitación motivó el desarrollo posterior de las redes LSTM.

\subsection{Redes Long Short-Term Memory (LSTM)}

Para mitigar las limitaciones de las RNN simples, se introdujeron las redes LSTM. Estas arquitecturas poseen bloques de celdas más complejos que les permiten retener más entradas pasadas.
 Una LSTM gestiona el flujo de información a través de tres compuertas principales:
\begin{enumerate}
    \item \textbf{Compuerta de Olvido (Forget Gate):} Decide qué información del estado de celda anterior ($C_{t-1}$) se conserva. Utiliza una función sigmoide que genera valores entre 0 (olvidar)
    y 1 (retener) basándose en la entrada actual $x_t$ y el estado oculto previo $h_{t-1}$.
    
    \item \textbf{Compuerta de Entrada (Input Gate):} Determina qué nueva información proveniente de $x_t$ se almacenará en el estado de la celda $C_t$. Combina una capa sigmoide (que decide
    qué actualizar) y una capa `tanh` (que genera los candidatos a añadir).
    
    \item \textbf{Compuerta de Salida (Output Gate):} Controla qué partes del estado interno de la celda ($C_t$) se expondrán como salida final ($h_t$). Primero se filtra $C_t$ con una tangente
    hiperbólica (para poner los valores entre -1 y 1) y luego se multiplica por la salida de la compuerta sigmoide, generando el estado oculto $h_t$ que pasará al siguiente paso de tiempo.
\end{enumerate}

A pesar de sus mejoras en la gestión del contexto, las LSTM y las RNN comparten una debilidad estructural inherente: su naturaleza secuencial impide la paralelización dentro de los ejemplos
de entrenamiento. El cálculo del estado oculto $h_t$ depende del estado anterior $h_{t-1}$, lo que hace que el tiempo de entrenamiento sea prohibitivo para secuencias muy largas y limite el 
uso eficiente del hardware moderno como las GPUs

\section{Arquitectura Transformer}

En 2017, \textcite{vaswani_attention_2023} propusieron un cambio de paradigma radical con la arquitectura \textbf{Transformer}. A diferencia de las RNN y LSTM, que procesan la información de
manera recurrente, el Transformer se basa únicamente en mecanismos de atención para establecer dependencias globales entre la entrada y la salida, prescindiendo por completo de la recurrencia
y las convoluciones \parencite{vaswani_attention_2023}.

\subsection{Superando la Secuencialidad: Paralelismo y Posición}

La limitación fundamental de las arquitecturas previas era su naturaleza secuencial: para calcular el estado oculto en el tiempo $t$, se requieria el estado en $t-1$. Esto impedía la 
paralelización durante el entrenamiento. El Transformer rompe esta restricción procesando todos los tokens de la secuencia simultáneamente. Esto reduce el número de operaciones secuenciales de
$O(n)$ (en RNNs) a un tiempo constante $O(1)$, permitiendo un entrenamiento significativamente más rápido y eficiente en hardware moderno (GPUs/TPUs) \parencite{vaswani_attention_2023}.

Sin embargo, al eliminar la recurrencia, el modelo pierde inherentemente la noción del orden de las palabras (es decir, procesa ''Gato come ratón'' igual que ''Ratón come gato''). Para solucionar
esto, el Transformer inyecta información sobre la posición relativa o absoluta de los tokens mediante \textbf{Codificaciones Posicionales} (\textit{Positional Encodings}).

Estas codificaciones se suman a los embeddings de entrada y, en la propuesta original, utilizan funciones sinusoidales de diferentes frecuencias:

\begin{equation}
    PE_{(pos, 2i)} = \sin(pos / 10000^{2i/d_{model}})
\end{equation}
\begin{equation}
    PE_{(pos, 2i+1)} = \cos(pos / 10000^{2i/d_{model}})
\end{equation}

Donde $pos$ es la posición y $i$ la dimensión. Esta elección permite al modelo aprender a atender a posiciones relativas fácilmente, ya que cualquier desplazamiento fijo puede representarse como
una función lineal \parencite{vaswani_attention_2023}.

\subsection{Mecanismo de Atención (Scaled Dot-Product Attention)}

El núcleo de la arquitectura es el mecanismo de atención, específicamente la \textit{Scaled Dot-Product Attention}. Este mecanismo permite al modelo relacionar diferentes posiciones de una 
secuencia para calcular una representación de la misma. 

La función de atención se describe como un mapeo de una consulta ($Q$, \textit{Query}) y un conjunto de pares clave-valor ($K$, \textit{Key}; $V$, \textit{Value}) a una salida. Matemáticamente:

\begin{equation}
    Attention(Q, K, V) = \text{softmax}\left(\frac{QK^T}{\sqrt{d_k}}\right)V
\end{equation}

Donde:
\begin{itemize}
    \item $QK^T$: Representa el producto punto entre la consulta y las claves, midiendo la similitud entre ellas.
    \item $\sqrt{d_k}$: Es un factor de escalado (donde $d_k$ es la dimensión de las claves). Este término es crucial porque, para valores grandes de $d_k$, los productos punto pueden crecer en
    magnitud, empujando a la función softmax a regiones con gradientes extremadamente pequeños. El escalado contrarresta este efecto \parencite{vaswani_attention_2023}.
    \item \text{softmax}: Normaliza los pesos para que sumen 1.
\end{itemize}

Este cálculo permite al modelo ponderar la importancia de cada palabra en relación con todas las demás, independientemente de su distancia física en la secuencia, resolviendo efectivamente el 
problema de las dependencias a largo plazo.

\subsection{Atención Multi-Cabeza (Multi-Head Attention)}

En lugar de aplicar una única función de atención con claves, valores y consultas de dimensión completa ($d_{model}$), \textcite{vaswani_attention_2023} demostraron que es beneficioso proyectar
linealmente estas entradas $h$ veces con diferentes proyecciones aprendidas. Este mecanismo se denomina \textbf{Multi-Head Attention}.

La intuición detrás de este enfoque es permitir al modelo ''atender'' conjuntamente a información proveniente de diferentes subespacios de representación en diferentes posiciones. Con una sola
cabeza de atención, el promedio ponderado podría diluir esta información. Al tener múltiples cabezas, cada una puede especializarse en capturar distintos tipos de relaciones (e.g., una cabeza 
puede enfocarse en la concordancia gramatical sujeto-verbo, mientras que otra captura dependencias semánticas o anafóricas).

Matemáticamente, dado un número de cabezas $h$, el proceso se descompone en tres etapas:

\begin{enumerate}
    \item \textbf{Proyecciones Lineales:} Las consultas ($Q$), claves ($K$) y valores ($V$) se proyectan a dimensiones menores ($d_k, d_v$) utilizando matrices de pesos entrenables específicas
    para cada cabeza $i$:
    \begin{equation}
        Q_i = QW_i^Q, \quad K_i = KW_i^K, \quad V_i = VW_i^V
    \end{equation}
    Donde las matrices de proyección son $W_i^Q \in \mathbb{R}^{d_{model} \times d_k}$, $W_i^K \in \mathbb{R}^{d_{model} \times d_k}$ y $W_i^V \in \mathbb{R}^{d_{model} \times d_v}$.

    \item \textbf{Atención Escalada en Paralelo:} Se aplica la función de atención independientemente a cada conjunto proyectado para obtener la salida de cada cabeza ($head_i$):
    \begin{equation}
        head_i = \text{Attention}(Q_i, K_i, V_i)
    \end{equation}

    \item \textbf{Concatenación y Proyección Final:} Finalmente, las salidas de todas las cabezas se concatenan y se proyectan una vez más para recuperar la dimensión original del modelo. Esta
    proyección final mezcla la información capturada por los distintos subespacios:
    \begin{equation}
        MultiHead(Q, K, V) = \text{Concat}(head_1, \dots, head_h)W^O
    \end{equation}
    Donde $W^O \in \mathbb{R}^{hd_v \times d_{model}}$ es la matriz de pesos de salida.
\end{enumerate}

En la implementación original, se utilizaron $h=8$ cabezas y una dimensión reducida para cada una de $d_k = d_v = d_{model}/h = 64$. Debido a esta reducción de dimensión, el costo computacional
total es similar al de una atención de cabeza única con dimensión completa \parencite{vaswani_attention_2023}.
\subsection{Arquitectura Encoder-Decoder}

El Transformer original adopta una estructura de codificador-decodificador:

\begin{itemize}
    \item \textbf{Encoder:} Compuesto por una pila de $N=6$ capas idénticas. Cada capa tiene dos subcapas: un mecanismo de atención multi-cabeza y una red \textit{feed-forward} conectada por posición.
    Se emplean conexiones residuales y normalización de capa (\textit{LayerNorm}) alrededor de cada subcapa.
    
    \item \textbf{Decoder:} También compuesto por $N=6$ capas. Además de las dos subcapas del encoder, inserta una tercera subcapa que realiza atención multi-cabeza sobre la salida del encoder.
    Crucialmente, la subcapa de auto-atención en el decoder está \textbf{enmascarada} (\textit{Masked Multi-Head Attention}). Esto asegura que, para predecir la posición $i$, el modelo solo pueda 
    atender a las posiciones anteriores a $i$, preservando la propiedad auto-regresiva necesaria para la generación de texto \parencite{vaswani_attention_2023}.
\end{itemize}

Aunque el paper original presentó esta estructura dual, la evolución posterior de los LLMs (como la serie GPT) ha tendido a utilizar arquitecturas \textit{Decoder-only}, optimizadas para tareas de 
generación, mientras que modelos como BERT utilizan arquitecturas \textit{Encoder-only} para tareas de comprensión.

\section{Grandes Modelos de Lenguaje: LLMs}
Los Grandes Modelos de Lenguaje (LLMs) representan la culminación actual de la arquitectura Transformer aplicada a escala masiva. Se definen como modelos probabilísticos entrenados en vastos corpus
de texto que aprenden la distribución de probabilidad sobre secuencias de palabras \parencite{atkinson2025}. A diferencia de sus predecesores, los LLMs exhiben capacidades emergentes—habilidades que
no fueron explícitamente programadas, como el razonamiento aritmético o la traducción que surgen únicamente como función de la escala del modelo y de los datos.

Aunque todos se basan en el Transformer, la elección de componentes define su naturaleza funcional:

\begin{itemize}
    \item \textbf{Arquitecturas Encoder-only (e.g., BERT):} Diseñadas para tareas de comprensión y clasificación. \textcite{devlin_bert_nodate} introdujeron el objetivo de 
    \textit{Masked Language Modeling} (MLM), donde el modelo tiene acceso bidireccional a todo el contexto para predecir tokens ocultos aleatoriamente (e.g., $P(w_i | w_1, ..., w_{i-1}, w_{i+1}, ... w_n)$).
    Esto permite representaciones semánticas profundas pero dificulta la generación de texto fluido.
    
    \item \textbf{Arquitecturas Decoder-only (e.g., GPT):} Optimizadas para generación. Utilizan una máscara causal (causal masking) que impide que la atención observe tokens futuros, forzando al modelo
    a aprender una tarea puramente autorregresiva ($P(w_t | w_{<t})$). Aunque teóricamente ven menos contexto por paso que BERT, su capacidad generativa los ha convertido en el estándar dominante para 
    LLMs modernos.
\end{itemize}


\subsection{Leyes de Escalado (Scaling Laws)}

El diseño eficiente de estos modelos se rige por leyes empíricas de potencia. Inicialmente, \textcite{kaplan_scaling_2020} observaron que la pérdida ($L$) del modelo disminuye predeciblemente como una ley de
potencia respecto al número de parámetros ($N$), el tamaño del dataset ($D$) y el cómputo ($C$), sugiriendo que aumentar el tamaño del modelo ($N$) era el factor más crítico para mejorar el rendimiento.

Sin embargo, \textcite{hoffmann_training_2022} refutaron parcialmente esta visión con el modelo \textbf{Chinchilla}. Demostraron que la mayoría de los LLMs de la época (como GPT-3) estaban ''sub-entrenados''.
Su hallazgo fundamental establece que, para un presupuesto computacional fijo $C$, el número de parámetros $N$ y la cantidad de datos $D$ deben escalarse en proporciones iguales (linealmente). Esto redefinió 
la frontera de eficiencia, indicando que los modelos más pequeños entrenados con muchos más datos son computacionalmente más óptimos que modelos gigantes con pocos datos.


\subsection{Ciclo de Vida de un LLM}

El desarrollo de un LLM moderno sigue un proceso estándar de tres etapas:

\subsubsection{Pre-Entrenamiento}
Esta es la fase fundamental y de mayor consumo computacional. El modelo se inicializa con pesos aleatorios y se entrena sobre corpus masivos (terabytes de texto provenientes de Common Crawl, Wikipedia, libros, código, etc.) para aprender la estructura estadística del lenguaje.

Dependiendo de la arquitectura, se optimizan funciones de pérdida distintas:
\begin{itemize}
    \item \textbf{Masked Language Modeling (MLM):} Utilizado por modelos \textit{Encoder-only} como BERT. El objetivo es reconstruir tokens ocultos en la entrada (e.g., $w_i$) basándose en el contexto bidireccional ($w_{1}, ..., w_{i-1}, w_{i+1}, ..., w_n$). Esto permite aprender representaciones 
    profundas de la sintaxis y semántica, pero no es ideal para generar texto \parencite{devlin_bert_nodate}.
    \item \textbf{Causal Language Modeling (CLM):} Utilizado por modelos \textit{Decoder-only} como GPT. El modelo intenta predecir el siguiente token $w_t$ dado únicamente el historial previo ($w_{<t}$), maximizando la probabilidad logarítmica negativa (\textit{negative log-likelihood}) sobre todo el
    dataset. Este enfoque fuerza al modelo a aprender razonamiento y coherencia narrativa.
\end{itemize}

El resultado es un ''modelo base'' que posee vastos conocimientos latentes pero que no es capaz de seguir instrucciones (e.g., ante la entrada ''¿Cómo hacer una torta?'', un modelo base podría completar con ''¿Cómo hacer una pizza?'', siguiendo la distribución probable de una lista de preguntas).
\subsubsection{Ajuste Fino (Fine-Tuning)}
Para transformar un modelo base en un asistente útil, se aplica el paradigma de \textit{Instruction Tuning}. A diferencia del aprendizaje por transferencia clásico (donde se entrenaba una nueva capa para una tarea única como clasificación), aquí se busca que el modelo generalice la capacidad de
seguir instrucciones arbitrarias. El método estándar es el \textit{Supervised Fine-Tuning} (SFT). Se utiliza un dataset curado de demostraciones humanas de alta calidad en formato (Prompt, Respuesta Ideal). El modelo ajusta sus pesos mediante descenso de gradiente para minimizar la diferencia entre 
su generación y la demostración humana. Según \textcite{ouyang_training_nodate}, esta etapa es crucial para definir el ''comportamiento'' y el estilo del modelo, aunque se utilizan muchos menos datos que en el pre-entrenamiento debido al requerimiento de intervención humana.

\subsubsection{Alineación y RLHF}
El SFT tiene limitaciones: escribir demostraciones perfectas es costoso y difícil. Para solucionar esto y alinear el modelo con la intención humana (seguridad, honestidad, utilidad), se emplea 
\textit{Reinforcement Learning from Human Feedback} (RLHF), proceso detallado por \textcite{ouyang_training_nodate} en tres pasos:
\begin{enumerate}
    \item \textbf{Recolección de Preferencias:} Los humanos comparan múltiples respuestas del modelo para una misma pregunta y las clasifican de mejor a peor.
    \item \textbf{Modelo de Recompensa (Reward Model - RM):} Se entrena una red neuronal separada para predecir la preferencia humana (es decir, asignar un escalar de recompensa $r$ a una 
    generación dada).
    \item \textbf{Optimización con PPO:} Se utiliza el algoritmo \textit{Proximal Policy Optimization} (PPO) para ajustar el LLM. El objetivo es maximizar la recompensa del RM sin desviarse 
    excesivamente de la distribución del modelo original (penalización KL), evitando así el ''reward hacking'' o la degradación del lenguaje.
\end{enumerate}


\subsection{Paradigmas de Inferencia (In-Context Learning)}

Una propiedad distintiva de los LLMs es su capacidad de \textit{In-Context Learning}: pueden adaptarse a nuevas tareas en tiempo de inferencia sin actualizar sus pesos, simplemente condicionándose a los
ejemplos provistos en el prompt \parencite{brown_language_nodate}.

\begin{itemize}
    \item \textbf{Zero-shot:} El modelo realiza la tarea solo con una descripción en lenguaje natural, sin ejemplos previos.
    \item \textbf{One-shot:} Se proporciona un único ejemplo de entrada-salida en el contexto.
    \item \textbf{Few-shot:} Se proporcionan varios ejemplos (generalmente entre 10 y 100) para guiar la generación del modelo.
\end{itemize}

\subsection{Técnicas Avanzadas y Variantes}

\subsubsection{RAG (Retrieval-Augmented Generation)}
Los LLMs estándar poseen únicamente \textit{memoria paramétrica}: todo su conocimiento está codificado estáticamente en sus pesos durante el entrenamiento. Esto los hace propensos a ''alucinaciones''
y obsolescencia.

\textcite{lewis_retrieval-augmented_nodate} formalizaron la arquitectura RAG para combinar modelos generativos con una \textit{memoria no paramétrica} externa (generalmente un índice vectorial denso de documentos,
como bases de datos clínicas). El proceso de inferencia cambia radicalmente:
\begin{enumerate}
    \item \textbf{Recuperación:} Dada una consulta $x$, un modelo \textit{Retriever} (e.g., DPR) busca en el índice los $k$ documentos más relevantes $z$.
    \item \textbf{Aumentación:} Estos documentos se concatenan con la consulta original.
    \item \textbf{Generación:} El LLM generador produce la respuesta $y$ condicionada tanto en la consulta como en la evidencia recuperada, aproximando $P(y|x) \approx \sum_{z} P(y|x,z)P(z|x)$.
\end{enumerate}
Esto permite actualizar el conocimiento del modelo simplemente actualizando la base de datos externa, sin necesidad de re-entrenar la red neuronal \parencite{lewis_retrieval-augmented_nodate}.

\subsubsection{Small Language Models (SLMs)}
Recientemente, ha surgido interés en modelos compactos (e.g., < 3B parámetros) que puedan ejecutarse en dispositivos locales. Trabajos como el de \textcite{gunasekar_textbooks_2023} con la serie
\textbf{Phi} han demostrado que la calidad de los datos ("textbook quality") es más determinante que la cantidad. Un modelo pequeño entrenado con datos sintéticos altamente curados puede superar a modelos
mucho más grandes entrenados con datos web ruidosos.

\subsection{Métricas de Evaluación}

La evaluación de modelos generativos presenta desafíos únicos respecto a tareas deterministas (como la clasificación), dado que no existe una única respuesta correcta. Para cuantificar el desempeño,
se emplean métricas automáticas que miden la probabilidad, la precisión y la cobertura (recall).

\subsubsection{Perplejidad (Perplexity - PPL)}
Es la métrica estándar intrínseca para evaluar Modelos de Lenguaje (LMs). Cuantifica la incertidumbre del modelo al predecir el siguiente token en una secuencia real. Matemáticamente, se define como
la exponenciación de la entropía cruzada promedio por token.

Dada una secuencia de prueba $X = (x_1, x_2, \dots, x_N)$, la perplejidad se calcula como:

\begin{equation}
    PPL(X) = \exp\left( -\frac{1}{N} \sum_{i=1}^{N} \ln P(x_i | x_{<i}) \right)
\end{equation}

Donde $P(x_i | x_{<i})$ es la probabilidad asignada por el modelo al token correcto $x_i$ dado el contexto previo.
\begin{itemize}
    \item Un valor bajo indica que el modelo asigna alta probabilidad a la secuencia de texto real (está ''menos perplejo'').
    \item Si el modelo fuera completamente aleatorio (asignando probabilidad uniforme $1/|V|$ a cada palabra del vocabulario), la perplejidad sería igual al tamaño del vocabulario $|V|$.
\end{itemize}

\subsubsection{BLEU (Bilingual Evaluation Understudy)}
Originalmente propuesta para traducción automática por \textcite{papineni_bleu_2001}, BLEU mide la \textbf{precisión} de n-gramas del texto generado comparándolo con una o más referencias humanas.

El cálculo penaliza la repetición y la longitud corta. La fórmula general es:

\begin{equation}
    BLEU = BP \cdot \exp\left( \sum_{n=1}^{N} w_n \log p_n \right)
\end{equation}

Donde:
\begin{itemize}
    \item $p_n$: Precisión de n-gramas (proporción de n-gramas generados que aparecen en la referencia).
    \item $w_n$: Peso asignado a cada tamaño de n-grama (usualmente uniforme $1/N$).
    \item $BP$ (Brevity Penalty): Factor de penalización si el texto generado ($c$) es más corto que la referencia ($r$):
    \begin{equation}
        BP = \begin{cases} 
        1 & \text{si } c > r \\
        e^{(1 - r/c)} & \text{si } c \le r 
        \end{cases}
    \end{equation}
\end{itemize}
BLEU es estándar de facto pero ha sido criticada por no capturar significado semántico, enfocándose solo en coincidencias léxicas exactas.

\subsubsection{ROUGE (Recall-Oriented Understudy for Gisting Evaluation)}
A diferencia de BLEU, ROUGE se enfoca en la \textbf{cobertura (recall)}, siendo la métrica predominante en tareas de resumen automático \parencite{lin_rouge_nodate}. Mide cuánto del contenido de
la referencia humana fue capaz de recuperar el modelo.

La variante más común, ROUGE-N, se calcula como:

\begin{equation}
    \text{ROUGE-N} = \frac{\sum_{S \in \{Referencias\}} \sum_{gram_n \in S} Count_{match}(gram_n)}{\sum_{S \in \{Referencias\}} \sum_{gram_n \in S} Count(gram_n)}
\end{equation}

Donde:
\begin{itemize}
    \item El numerador cuenta los n-gramas coincidentes entre la generación y la referencia.
    \item El denominador cuenta el total de n-gramas en la \textbf{referencia} (a diferencia de BLEU, que divide por los generados).
\end{itemize}
Existen variantes como \textbf{ROUGE-L}, que utiliza la Subsecuencia Común Más Larga (LCS) para capturar la estructura de la oración sin depender de n-gramas fijos.


\section{Modelos de Visión-Lenguaje (VLLMs)}
La convergencia entre la visión por computadora y el procesamiento del lenguaje natural ha dado lugar a los Modelos de Visión-Lenguaje (VLLMs). Estos sistemas buscan trascender las tareas visuales
aisladas (como clasificar una imagen en una de 1000 categorías fijas) para lograr una comprensión semántica abierta, donde las imágenes se entienden en el contexto del lenguaje humano.

\subsection{El Transformer Visual (ViT)}

Históricamente, la visión por computadora estuvo dominada por las Redes Neuronales Convolucionales (CNNs), las cuales procesan imágenes mediante filtros locales que asumen una invariancia espacial.
Sin embargo, tras el éxito rotundo del Transformer en NLP, \textcite{dosovitskiy_image_2021} cuestionaron si esta dependencia de las CNNs era necesaria.

En su trabajo ''\textit{An Image is Worth 16x16 Words}'', propusieron el \textbf{Vision Transformer (ViT)}. La idea central consiste en tratar una imagen no como una cuadrícula de píxeles, sino
como una secuencia de ''palabras'' visuales. Para ello:
\begin{enumerate}
    \item La imagen se divide en parches fijos (e.g., $16 \times 16$ píxeles).
    \item Cada parche se aplana y se proyecta linealmente a un embedding (similar a los embeddings de palabras).
    \item Se añade información posicional (ya que el Transformer no tiene noción inherente de espacio 2D).
    \item La secuencia resultante se procesa con un Encoder de Transformer estándar.
\end{enumerate}

Los autores demostraron que, cuando se entrena a escala masiva (e.g., en JFT-300M), ViT supera a las arquitecturas convolucionales de estado del arte (como ResNet), con un costo computacional menor
\parencite{dosovitskiy_image_2021}. Este avance fue fundamental porque unificó la arquitectura base para visión y texto, facilitando su integración posterior.

\subsection{Alineación Multimodal: CLIP}

Aunque ViT permitió usar Transformers en visión, los modelos seguían entrenándose para predecir un conjunto cerrado de etiquetas (clasificación supervisada). \textcite{radford_learning_nodate}
revolucionaron este paradigma con \textbf{CLIP (Contrastive Language-Image Pre-training)}.

La hipótesis de CLIP es que la supervisión del lenguaje natural es mucho más rica y escalable que las etiquetas manuales. El modelo se entrena con un objetivo \textit{contrastivo} sobre un 
dataset masivo de 400 millones de pares (imagen, texto) recolectados de internet.
El funcionamiento es el siguiente:
\begin{itemize}
    \item Un \textbf{Encoder de Imagen} (que puede ser un ViT o ResNet) convierte la imagen en un vector.
    \item Un \textbf{Encoder de Texto} (Transformer) convierte la descripción en otro vector.
    \item El objetivo es maximizar la similitud coseno entre los vectores de pares correctos (imagen, su texto real) y minimizarla para los pares incorrectos en un lote de $N$ ejemplos.
\end{itemize}

Esto dota al modelo de capacidades \textit{Zero-shot} asombrosas: CLIP puede clasificar imágenes en categorías que nunca vio durante el entrenamiento simplemente proporcionándole los nombres de
las clases como texto (e.g., ''una foto de un [melanoma]''), alineando efectivamente el espacio visual con el semántico \parencite{radford_learning_nodate}.

\subsection{Grandes Modelos Multimodales (LMMs) y Visual Instruction Tuning}

El último salto evolutivo consiste en dotar a los LLMs de ''ojos''. Si bien CLIP puede alinear imágenes y texto, no puede mantener una conversación compleja sobre una imagen. Aquí surgen modelos
como \textbf{LLaVA (Large Language and Vision Assistant)}, propuesto por \textcite{liu_visual_2023}.

LLaVA introduce el concepto de \textit{Visual Instruction Tuning}. La arquitectura conecta un encoder visual pre-entrenado (como CLIP ViT-L/14) con un LLM (como Vicuna/LLaMA) mediante una capa
de proyección lineal simple. Esta proyección transforma las características visuales al mismo espacio de embeddings que las palabras, permitiendo que la imagen actúe como un ''prompt visual'' 
para el LLM.

El entrenamiento se realiza en dos etapas:
\begin{enumerate}
    \item \textbf{Pre-entrenamiento de alineación:} Se congelan el encoder visual y el LLM, y se entrena solo la matriz de proyección para alinear las características visuales con el espacio 
    de palabras.
    \item \textbf{Ajuste Fino de Instrucciones (End-to-End):} Se entrena el modelo completo (o partes de él) con datos de instrucciones multimodales generados sintéticamente por GPT-4 
    (pares de preguntas y respuestas complejas sobre imágenes).
\end{enumerate}

Este enfoque permite que el modelo no solo reconozca objetos, sino que razone sobre ellos, responda preguntas abiertas y siga instrucciones complejas sobre el contenido visual, comportándose
como un verdadero chatbot multimodal \parencite{liu_visual_2023}.

\printbibliography

\end{document}
